% foundations/state_model.tex

\chapter{State Model}
\label{chap:state-model}

\index{State!model}
\index{Engineering state}

An alignment is not described by geometry alone. At each arc-length
position, the theory requires a local description from which geometric
orientation and subsequent spatial interpretation can be developed. The
state model provides this description without incorporating realization,
representation, runtime architecture, or optimization into the state
concept itself.

% ============================================================
\section{Engineering State}
% ============================================================

\begin{definition}[Engineering state]
\label{def:state}
\index{State!engineering}
An engineering state is a structured description of an engineering
object under a specified condition and within a specified context.
\end{definition}

A state is neither the engineering object itself nor a representation of
that object. It selects the quantities required to describe a particular
condition. Different questions may therefore require different state
spaces while referring to the same engineering object.

For alignment geometry, the independent variable is intrinsic
arc-length. A state trajectory is consequently a mapping
\[
x:I\rightarrow\mathcal X,
\qquad
s\longmapsto x(s),
\]
where \(\mathcal X\) is the state space appropriate to the geometric or
engineering question under consideration.

% ============================================================
\section{Planar Alignment State}
% ============================================================

\begin{definition}[pose2]
\label{def:notation-pose2}
\index{State!pose2}
The planar alignment state is
\begin{equation}
\mathrm{pose2}(s)
=
\bigl(\gamma(s),\theta(s)\bigr),
\label{eq:pose2-definition}
\end{equation}
where \(\gamma(s)\in\mathbb{R}^{2}\) is the planar position and
\(\theta(s)\) the heading.
\end{definition}

The evolution of \(\mathrm{pose2}\) follows directly from the axiomatic
foundation:
\begin{equation}
\begin{aligned}
\gamma'(s)
&=
\begin{pmatrix}
\cos\theta(s)\\
\sin\theta(s)
\end{pmatrix},\\
\theta'(s)
&=\kappa(s).
\end{aligned}
\label{eq:pose2-evolution}
\end{equation}

Curvature is therefore the control quantity governing the directional
evolution of the planar alignment state. Given an initial pose and a
curvature function, the complete planar state trajectory is determined.

% ============================================================
\section{Spatial Railway State}
% ============================================================

A spatial railway state requires more than planar position and heading.
Profile, cant, and the applicable realization context contribute to the
construction of a spatial position and local frame. These dependencies
are developed in detail in the dedicated state chapters. At foundation
level, the spatial extension is expressed by adding the cant rotation to
the planar state.

\begin{definition}[Foundational pose3 state]
\label{def:pose3}
\index{State!pose3}
The foundational spatial railway state is
\begin{equation}
\mathrm{pose3}(s)
=
\bigl(\gamma(s),\theta(s),\phi(s)\bigr),
\label{eq:pose3-definition}
\end{equation}
where \(\phi(s)\) denotes cant rotation about the local tangent.
\end{definition}

This tuple records the intrinsic quantities required for the subsequent
spatial construction. It is not yet a complete world-space pose. The
spatial position, local frame, and physical interpretation arise only
when profile and realization context are applied.

% ============================================================
\section{State Controls}
% ============================================================

\begin{definition}[Foundational controls]
\label{def:control-quantities}
\index{Control quantities!foundational}
\index{Curvature!as control quantity}
The foundational control pair is
\begin{equation}
u(s)=\bigl(\kappa(s),\omega(s)\bigr),
\qquad
\omega(s)=\phi'(s).
\label{eq:pose3-controls}
\end{equation}
\end{definition}

Curvature controls heading evolution, while cant-rate controls the
evolution of cant rotation. The corresponding foundational system is
\begin{equation}
\begin{aligned}
\gamma'(s)
&=
\begin{pmatrix}
\cos\theta(s)\\
\sin\theta(s)
\end{pmatrix},\\
\theta'(s)&=\kappa(s),\\
\phi'(s)&=\omega(s).
\end{aligned}
\label{eq:pose3-evolution}
\end{equation}

This formulation establishes an arc-length-driven state system. It does
not imply that every engineering quantity is a control or that every
spatial quantity belongs to the intrinsic state.

% ============================================================
\section{Intrinsic and Realized Descriptions}
% ============================================================

\begin{definition}[Intrinsic state description]
\index{State!intrinsic description}
An intrinsic state description uses quantities defined relative to the
alignment and does not require a particular world-space realization.
\end{definition}

Curvature, heading evolution, cant evolution, and the ordered relation of
states along arc-length are intrinsic. Their mathematical meaning can be
established before a coordinate reference system or physical embedding
is selected.

\begin{definition}[Realized state description]
\index{State!realized description}
A realized state description expresses an engineering state within a
specified metric realization context.
\end{definition}

A realized state may contain world-space position, a spatial frame, or
other metric quantities. Such quantities are related to the intrinsic
state through realization operators; they do not become intrinsic merely
because they are evaluated from it.

The distinction is deliberately one of description rather than of two
unrelated state theories. The same alignment may admit multiple realized
descriptions while retaining one intrinsic engineering identity.

% ============================================================
\section{State, Observation, and Representation}
% ============================================================

A state may be evaluated, observed, or represented. These operations
must remain conceptually distinct.
\index{Observation!state observation}
\index{Representation!state representation}

An evaluated state is derived from engineering knowledge and applicable
operators. An observed state records information obtained from a source
such as measurement or inspection. A representation communicates or
processes state information in a particular form. None of these forms is
identical merely because they share numerical values at a given instant.

This separation is essential for later treatment of measured reality,
physical realization, runtime evaluation, and optimization. Those
chapters may operate on or compare states, but they do not redefine the
foundational state concept established here.

% ============================================================
\section{Canonical Interpretation}
% ============================================================

\begin{proposition}
A planar railway alignment can be interpreted as an intrinsic state
trajectory parameterized by arc-length and controlled by curvature.
\end{proposition}
\index{State trajectory}

\begin{proposition}
Cant evolution extends the foundational state without by itself
establishing a complete metric or physical spatial realization.
\end{proposition}

\begin{proposition}
Realization, observation, representation, and optimization relate to
engineering states through explicit operators and relations; they are
not constituent parts of the foundational state definition.
\end{proposition}

The state model therefore provides a common mathematical interface
between intrinsic geometry and the later chapters on spatial
interpretation, realization, runtime evaluation, and engineering
assessment while preserving the responsibility boundaries established
by the Knowledge Kernel.
